\specialchapter{尼古拉·布尔巴基 数学原理}

已出版的英文译本:

一般拓扑学, 第一部分

一般拓扑学, 第二部分

集合论

交换代数

代数, 第一部分

李群与李代数, 第一部分

\textbf{一般拓扑学, 第一部分}

\textbf{第一章. 拓扑结构}

§ 1. 开集, 邻域, 闭集

§ 2. 连续函数

§ 3. 子空间, 商空间

§ 4. 拓扑空间的乘积

§ 5. 开映射和闭映射

§ 6. 滤子

§ 7. 极限

§ 8. Hausdorff 空间和正则空间

§ 9. 紧空间和局部紧空间

§ 10. 真映射

§ 11. 连通性

习题

\textbf{第二章. 一致结构}

§ 1. 一致空间

§ 2. 一致连续函数

§ 3. 完备空间

§ 4. 一致空间与紧空间之间的关系

\textbf{第三章. 拓扑群}

§ 1. 群上的拓扑

§ 2. 子群, 商群, 同态, 齐性空间, 乘积群

§ 3. 群上的一致结构

§ 4. 在拓扑空间上适当作用的群; 拓扑群和带算子空间中的紧性

§ 5. 交换群中的无穷和

§ 6. 带算子的拓扑群; 拓扑环, 除环和域

\textbf{§ 7. 拓扑群和拓扑环的逆极限}

\textbf{第四章. 实数}

§ 1. 实数的定义

§ 2. 实直线的基本拓扑性质

§ 3. 实数域

§ 4. 扩展实直线

§ 5. 实值函数

§ 6. 连续和半连续实值函数

§ 7. 实数的无穷和与无穷积

§ 8. 实数的通常展开; R 的幂

习题

第五章. 单参数群 § 1. R 的子群和商群 § 2. 量的测量 § 3. 群 R 和 T 的拓扑刻画 § 4. 指数和对数 习题

第六章. 实数空间和射影空间 § 1. 实数空间 R \(^{n}\) § 2. 欧氏距离, 球和球面 § 3. 实射影空间 习题

第七章. 加法群 R \(^{n}\) § 1. R \(^{n}\) 的子群和商群 § 2. R \(^{n}\) 及其商群的连续同态 § 3. 群 R \(^{n}\) 中的无穷和 习题

第八章. 复数 § 1. 复数, 四元数 § 2. 角度测量, 三角函数 § 3. 复数的无穷和与无穷积 § 4. 复数空间和射影空间 习题

第九章. 实数在一般拓扑学中的应用 § 1. 由一族伪度量生成一致性; 可一致化空间 § 2. 度量空间和可度量化空间 § 3. 可度量化群, 赋值域, 赋范空间和代数 § 4. 正规空间 § 5. Baire 空间 § 6. Polish 空间, Souslin 空间, Borel 集 附录: 赋范代数中的无穷积 习题

第十章. 函数空间 § 1. S-收敛的一致性 § 2. 等度连续集 § 3. 特殊函数空间 § 4. 连续实值函数的逼近 习题

\textbf{一般拓扑学, 第二部分}

\textbf{集合论}

\textbf{第一章. 形式数学的描述}

§ 1. 项和关系

§ 2. 定理

§ 3. 逻辑理论

§ 4. 量化理论

§ 5. 等式理论

附录. 项和关系的刻画 习题

第二章. 集合论

§ 1. 集合化关系

§ 2. 有序对

§ 3. 对应

§ 4. 一族集合的并与交

§ 5. 一族集合的乘积

§ 6. 等价关系

习题

第三章. 有序集, 基数, 整数

§ 1. 序关系. 有序集

§ 2. 良序集

§ 3. 等势集. 基数

§ 4. 自然整数. 有限集

§ 5. 整数的性质

§ 6. 无限集

§ 7. 逆极限和直极限

习题

第四章. 结构

§ 1. 结构和同构

§ 2. 态射和导出结构

§ 3. 泛映射

习题

结果概要

引言

§ 1. 集合的元素和子集

§ 2. 函数

§ 3. 集合的乘积

§ 4. 一族集合的并、交、乘积

§ 5. 等价关系和商集

§ 6. 有序集

§ 7. 幂. 可数集

§ 8. 集合的尺度. 结构

\textbf{交换代数}

第一章. 平坦模 § 1. 图表和正合序列 § 2. 平坦模 § 3. 忠实平坦模 § 4. 平坦模与 ``Tor'' 函子 习题

第二章. 局部化 § 1. 素理想 § 2. 分式环和分式模 § 3. 局部环. 从局部到整体的过渡 § 4. 环的谱和模的支集 § 5. 有限生成投射模. 可逆分式理想 习题

第三章. 分次、滤过和拓扑 § 1. 有限生成分次代数 § 2. 滤过环和模的一般结果 § 3. 诺特环上的 m-进拓扑 § 4. 完备环中的提升 § 5. 滤过模的平坦性 习题

第四章. 关联素理想与初等分解 § 1. 与模相关的素理想 § 2. 初等分解 § 3. 分次模中的初等分解 习题

第五章. 整数 § 1. 整元的概念 § 2. 素理想的提升 § 3. 域上的有限生成代数 习题

第六章. 赋值 § 1. 赋值环 § 2. 位 § 3. 赋值 § 4. 赋值的高度 § 5. 由赋值定义的拓扑 § 6. 绝对值 § 7. 逼近定理 § 8. 赋值到代数扩张的延拓 § 9. 应用: 局部紧域 § 10. 赋值到超越扩张的延拓 习题

\textbf{交换代数}

第七章. 除子

§ 1. Krull 整环

§ 2. Dedekind 整环

§ 3. 唯一分解整环

§ 4. 整闭诺特整环上的模 习题

\textbf{代数, 第一部分}

\textbf{第一章. 代数结构}

§ 1. 复合律; 结合律; 交换律

§ 2. 单位元; 可消去元; 可逆元

§ 3. 作用

§ 4. 群和带算子的群

§ 5. 在集合上作用的群

§ 6. 扩张, 可解群, 幂零群

§ 7. 自由幺半群, 自由群

§ 8. 环

§ 9. 域

§ 10. 逆极限和直极限

习题

\textbf{第二章. 线性代数}

§ 1. 模

§ 2. 线性映射的模. 对偶性

§ 3. 张量积

§ 4. 张量积与同态模之间的关系

§ 5. 标量环的扩张

§ 6. 模的逆极限和直极限

§ 7. 向量空间

§ 8. 向量空间中标量域的限制

§ 9. 仿射空间和射影空间

§ 10. 矩阵

§ 11. 分次模和环

附录. 伪模

习题

\textbf{第三章. 张量代数、外代数、对称代数}

§ 1. 代数

§ 2. 代数的例子

§ 3. 分次代数

§ 4. 代数的张量积

§ 5. 张量代数. 张量

§ 6. 对称代数

§ 7. 外代数

§ 8. 行列式

§ 9. 范数和迹

§ 10. 导子

§ 11. 余代数, 多重线性型的乘积, 内积和对偶性 附录. 交代数. 八元数

习题
